\chapter{Tiny210 MT7601U USB Wi-Fi
驱动}\label{tiny210-mt7601u-usb-wi-fi-ux9a71ux52a8}

\section{当前支持范围}\label{ux5f53ux524dux652fux6301ux8303ux56f4}

Tiny210 实板构建已经把以下组件连接成一条完整链路：

\begin{lstlisting}
S5PV210 USB Host PHY1
  -> EHCI host controller（USB-A，0xEC200000，IRQ 55）
  -> USB core / usb_bus
  -> MT7601U USB driver
  -> SoftMAC 扫描、Open-System 认证、关联
  -> WPA2-PSK/CCMP 四次握手
  -> net_device "wlan0"
  -> Ethernet/ARP/IPv4/UDP/ICMP/TCP
  -> DHCP、ping 和其他 socket 用户程序
\end{lstlisting}

默认目标是高速 USB MT7601U。QEMU \passthrough{\lstinline!cubieboard!}
替身没有这套 USB host， \passthrough{\lstinline!tiny210-qemu!} 不会验证
Wi-Fi；枚举、固件、射频和真实收发必须在 Tiny210 实板测试。

主要源码：

\begin{itemize}
\tightlist
\item
  \href{../../../kernel/arch/arm/mach-s5pv210/devices.c}{\passthrough{\lstinline!devices.c!}}：PHY1、EHCI
  resource。
\item
  \href{../../../kernel/drivers/usb/host/ehci-hcd.c}{\passthrough{\lstinline!ehci-hcd.c!}}：EHCI
  枚举和 bulk 传输。
\item
  \href{../../../kernel/drivers/usb/core/usb.c}{\passthrough{\lstinline!usb.c!}}：USB
  bus、driver 匹配和 URB 生命周期。
\item
  \href{../../../kernel/drivers/net/wireless/mt7601u.c}{\passthrough{\lstinline!mt7601u.c!}}：芯片、SoftMAC、WPA2
  和 netdev。
\item
  \href{../../../kernel/drivers/net/wireless/wpa_crypto.c}{\passthrough{\lstinline!wpa\_crypto.c!}}：PBKDF2/PTK/MIC/AES
  unwrap。
\item
  \href{../../../kernel/drivers/base/firmware.c}{\passthrough{\lstinline!firmware.c!}}：\passthrough{\lstinline!request\_firmware()!}。
\item
  \href{../../../kernel/net.c}{\passthrough{\lstinline!net.c!}}：net\_device、NAPI、ARP/IP/DHCP
  和 socket 网络栈。
\item
  \href{../../../user/wifi.c}{\passthrough{\lstinline!wifi.c!}}：用户态连接命令。
\end{itemize}

\section{硬件与 USB 拓扑}\label{ux786cux4ef6ux4e0e-usb-ux62d3ux6251}

默认使用 Tiny210 底板 USB-A 口：

\begin{figure}[htbp]
\centering
\resizebox{0.90\linewidth}{!}{\begin{tikzpicture}[node distance=8mm]
  \node[hbox] (cpu) {Cortex-A8};
  \node[box,right=of cpu] (ehci) {S5PV210 EHCI\\\texttt{0xEC200000}, IRQ 55};
  \node[box,right=of ehci] (hub) {可选 High-Speed Hub};
  \node[box,right=of hub] (mt) {MT7601U USB\\bulk IN/OUT + EP0};
  \node[hbox,right=of mt] (rf) {2.4 GHz\\802.11b/g/n};
  \node[gbox,above=5mm of ehci] (phy) {USB Host PHY1\\24 MHz reference};
  \draw[arr] (cpu)--(ehci); \draw[arr] (phy)--(ehci); \draw[arr] (ehci)--(hub);
  \draw[arr] (hub)--(mt); \draw[arr] (mt)--(rf);
\end{tikzpicture}}
\caption{Tiny210 USB-A 到 MT7601U 的硬件拓扑}
\label{fig:tiny-wifi-hw}
\end{figure}


\passthrough{\lstinline!host\_phy\_init()!} 会：

\begin{enumerate}
\def\labelenumi{\arabic{enumi}.}
\tightlist
\item
  打开 USBHOST 时钟 gate；
\item
  解除 PHY1 isolation；
\item
  选择 24 MHz reference；
\item
  清 PHY1 suspend/power-down；
\item
  pulse PHY1 和 host link reset。
\end{enumerate}

EHCI device resource 是：

{\def\LTcaptype{none} % do not increment counter
\begin{longtable}[]{@{}p{0.34\linewidth}p{0.60\linewidth}@{}}
\toprule\noalign{}
资源 & 值 \\
\midrule\noalign{}
\endhead
\bottomrule\noalign{}
\endlastfoot
EHCI MMIO & \passthrough{\lstinline!0xEC200000..0xEC2fffff!} \\
IRQ & 55（\passthrough{\lstinline!uhost!}） \\
USB PHY MMIO & \passthrough{\lstinline!0xEC100000..0xEC100fff!} \\
\end{longtable}
}

\section{USB/Wi-Fi 分层 Trace 开关}\label{tiny210-wifi-trace}

Tiny210 构建提供一个总开关和五个可独立覆盖的分层开关。一次打开 USB
探测、枚举、传输以及 MT7601U 配置和数据路径：

\begin{lstlisting}[language=sh]
make tiny210 USB_WIFI_TRACE=1
\end{lstlisting}

等级 2 进一步打开逐寄存器和逐 802.11 帧日志。若只编译镜像、不执行 SD
卡安装流程，可直接指定 Tiny210 Makefile 的镜像目标：

\begin{lstlisting}[language=sh]
make tiny210 USB_WIFI_TRACE=2
make -f ports/tiny210/Makefile USB_WIFI_TRACE=1 \
  build/tiny210/xv6-tiny210.uImage
\end{lstlisting}

\begingroup\small
\begin{longtable}{p{35mm}p{24mm}p{73mm}}
\toprule
开关 & 日志前缀 & 覆盖阶段 \\ \midrule \endhead
{\ttfamily\scriptsize USB\_CORE\_TRACE} & {\ttfamily\scriptsize usb-core:} & device/driver 注册、ID 匹配、probe/remove \\
{\ttfamily\scriptsize USB\_ENUM\_TRACE} & {\ttfamily\scriptsize usb-enum:} & EHCI 根口、Hub、描述符、地址和端点枚举 \\
{\ttfamily\scriptsize USB\_XFER\_TRACE} & {\ttfamily\scriptsize usb-xfer:} & control/bulk、QH/qTD、RX arm/completion 和 IRQ \\
{\ttfamily\scriptsize MT7601U\_\allowbreak CONFIG\_\allowbreak TRACE=1} & {\ttfamily\scriptsize wifi-config:} & 固件、eFuse、MAC/BBP/RF、校准和状态机 \\
{\ttfamily\scriptsize MT7601U\_\allowbreak CONFIG\_\allowbreak TRACE=2} & {\ttfamily\scriptsize wifi-reg:} & 在配置跟踪上增加逐寄存器访问 \\
{\ttfamily\scriptsize MT7601U\_\allowbreak DATA\_\allowbreak TRACE=1} & {\ttfamily\scriptsize wifi-data:} & TX/RX、NAPI、Ethernet 与 USB 缓冲流转 \\
{\ttfamily\scriptsize MT7601U\_\allowbreak DATA\_\allowbreak TRACE=2} & {\ttfamily\scriptsize wifi-frame:} & 在数据跟踪上增加逐 802.11 帧摘要 \\
\bottomrule
\end{longtable}
\endgroup

总开关只是五项的默认值，因此可以打开总体跟踪后单独关闭最密集的 EHCI
传输日志：

\begin{lstlisting}[language=sh]
make tiny210 USB_WIFI_TRACE=1 USB_XFER_TRACE=0
\end{lstlisting}

日志前缀与函数调用链的关系如下：

\begin{lstlisting}
ehci_enumerate_device
  -> usb-enum: descriptor/address/config/endpoint
  -> usb_device_register
     -> usb-core: match 148f:7601 -> mt7601u_probe
        -> wifi-config: firmware/eFuse/MAC/BBP/RF/calibration
           -> usb-xfer: EP0/EP8/EP4/EP5 QH and qTD

EHCI IRQ 55
  -> usb-xfer: EP4 completion
  -> mt7601u_rx_irq -> NAPI -> mt7601u_rx_parse
     -> wifi-data: Ethernet/NAPI path
     -> wifi-frame: each 802.11 frame (level 2 only)
\end{lstlisting}

错误、超时和初始化结果不受 Trace 开关影响。构建系统把当前五项取值写入
\passthrough{\lstinline!build/tiny210/.usb-wifi-trace!}；任一值变化都会自动重编译
\passthrough{\lstinline!usb.c!}、\passthrough{\lstinline!ehci-hcd.c!} 和
\passthrough{\lstinline!mt7601u.c!}，无需先执行
\passthrough{\lstinline!make clean!}。

根口上的 full/low-speed 设备需要 OHCI，当前没有 OHCI 驱动；高速 MT7601U
可直接由 EHCI 管理。也可以通过
\passthrough{\lstinline!S5PV210\_USB\_OTG=1!} 编入 OTG/DWC2 host
device，但默认 Wi-Fi 路径是 USB-A/EHCI。

\section{构建、firmware
和配置文件}\label{ux6784ux5efafirmware-ux548cux914dux7f6eux6587ux4ef6}

实板 \passthrough{\lstinline!ports/tiny210/Makefile!} 会编入：

\begin{lstlisting}
usb/core/usb
usb/host/ehci-hcd
usb/dwc2/hcd
net/wireless/mt7601u
net/wireless/wpa_crypto
base/firmware
net + workqueue + epoll + sysnet
wifi + dhcp + ping 用户程序
\end{lstlisting}

固件文件必须存在于仓库的：

\begin{lstlisting}
firmware/mt7601u.bin
\end{lstlisting}

构建会把它链接到内核，并在 \passthrough{\lstinline!make tiny210!}
时复制到已挂载的 FAT32 bootfs：

\begin{lstlisting}
/Volumes/bootfs/mt7601u.bin
\end{lstlisting}

\passthrough{\lstinline!request\_firmware("mt7601u.bin")!} 的优先级是：

\begin{lstlisting}
FAT32 根目录 /boot/mt7601u.bin
  -> 找不到或不可读
内核镜像中的 fw_mt7601u_start..fw_mt7601u_end
\end{lstlisting}

这样可以只替换 SD 卡上的固件进行测试，同时保留内置副本作为启动兜底。

可选 Wi-Fi 配置文件默认位于：

\begin{lstlisting}
ports/tiny210/wifi.conf
\end{lstlisting}

内容：

\begin{lstlisting}
ssid=your-access-point
psk=your-passphrase
\end{lstlisting}

该文件被 \passthrough{\lstinline!.gitignore!} 排除。构建发现它时，通过
\passthrough{\lstinline!mkfs --etc!} 写入根文件系统的
\passthrough{\lstinline!/etc/wifi.conf!}。不要把真实密码提交到仓库。

\section{initcall、枚举与驱动匹配}\label{initcallux679aux4e3eux4e0eux9a71ux52a8ux5339ux914d}

初始化顺序依靠 initcall level 和 order：

\begin{lstlisting}
netinit()                         注册 net bus
subsys_initcall("usbcore")       注册 usb bus
module_init_ord("mt7601u", 1)   注册 USB driver，建立 mt7601u_wq
module_init_ord("ehci-s5p", 2)  probe EHCI，申请 IRQ55，创建 ehci-enum
scheduler()                       CPU IRQ 打开，内核线程开始运行
ehci-enum                         PHY/host 初始化并枚举设备
\end{lstlisting}

EHCI 枚举流程：

\begin{figure}[htbp]
\centering
\resizebox{\linewidth}{!}{\begin{tikzpicture}[node distance=5mm]
  \node[hbox] (host) {EHCI host/PHY 初始化};
  \node[box,right=of host] (reset) {Port reset\\GET\_DESCRIPTOR(8)};
  \node[box,right=of reset] (addr) {SET\_ADDRESS\\读取 device/config};
  \node[box,right=of addr] (ep) {解析 bulk endpoint\\SET\_CONFIGURATION};
  \node[box,right=of ep] (bus) {\texttt{usb\_device\_register()}};
  \node[hbox,right=of bus] (probe) {VID:PID 匹配\\\texttt{mt7601u\_probe()}};
  \draw[arr] (host)--(reset); \draw[arr] (reset)--(addr); \draw[arr] (addr)--(ep);
  \draw[arr] (ep)--(bus); \draw[arr] (bus)--(probe);
\end{tikzpicture}}
\caption{EHCI 枚举 MT7601U 并交给 USB driver 的过程}
\label{fig:tiny-wifi-enum}
\end{figure}


MT7601U ID 表包含 \passthrough{\lstinline!148f:7601!} 以及多个 OEM
VID:PID。USB core 使用 ID table 匹配，不依赖设备名。EHCI
枚举在内核线程而不是 initcall 中完成，是因为端口上电、 reset
和描述符事务会等待数百毫秒，需要调度器、timer 和 IRQ 已可用。

\section{\texorpdfstring{\texttt{mt7601u\_probe()}
初始化阶段}{mt7601u\_probe() 初始化阶段}}\label{mt7601u_probe-ux521dux59cbux5316ux9636ux6bb5}

probe 的主要阶段如下：

\begin{figure}[htbp]
\centering
\begin{tikzpicture}[node distance=6mm]
  \node[hbox] (p) {\texttt{mt7601u\_probe}};
  \node[box,below=of p] (ver) {vendor control\\读取 ASIC/MAC version};
  \node[box,below=of ver] (fw) {校验并下载 firmware\\ILM + DLM + IVB};
  \node[box,below=of fw] (mcu) {等待 MCU running\\读取 256 B eFuse};
  \node[box,below=of mcu] (init) {WLAN PLL、MAC、USB DMA\\MAC/BBP/RF 表与校准};
  \node[hbox,below=of init] (rx) {channel 1，提交 4 个 RX buffer\\启动 delayed scan work};
  \draw[arr] (p)--(ver); \draw[arr] (ver)--(fw); \draw[arr] (fw)--(mcu);
  \draw[arr] (mcu)--(init); \draw[arr] (init)--(rx);
\end{tikzpicture}
\caption{MT7601U probe 的主要硬件初始化阶段}
\label{fig:tiny-wifi-probe}
\end{figure}


\subsection{控制传输}\label{ux63a7ux5236ux4f20ux8f93}

MT7601U 寄存器不是 CPU MMIO。驱动通过 USB vendor control request 读写：

\begin{lstlisting}
mt7601u_read32/write32
  -> udev->ops->control
  -> EHCI EP0 SETUP/DATA/STATUS qTD
  -> MT7601U vendor request
\end{lstlisting}

\subsection{固件下载}\label{ux56faux4ef6ux4e0bux8f7d}

固件头给出 ILM/DLM 长度。驱动把数据按最大
\passthrough{\lstinline!0x3800!} 字节分块，增加 FCE TXINFO， 通过
bulk-OUT 下载：

\begin{lstlisting}
ILM: 目的地址从 64 B IVB 之后开始
DLM: 目的地址 0x80000
最后: DEV_MODE=0x12 携带 64 B IVB，启动 MCU
\end{lstlisting}

每块之后检查 FCE/MCU 状态，最后轮询
\passthrough{\lstinline!MCU\_COM\_REG0 == 1!}。

\subsection{eFuse
和射频初始化}\label{efuse-ux548cux5c04ux9891ux521dux59cbux5316}

驱动读取 256 字节 eFuse，获得永久 MAC、country、晶振 frequency
offset、RSSI offset、参考温度和 LNA gain。随后初始化 MAC、BBP、RF
表，并通过 MCU command endpoint 执行 R、TXDCOC、LOFT、TXIQ、RXIQ 和 DPD
等校准。

任何阶段失败都会使 probe 返回失败，\passthrough{\lstinline!wlan0!}
不会出现。

\section{扫描、认证和 WPA2
状态机}\label{ux626bux63cfux8ba4ux8bc1ux548c-wpa2-ux72b6ux6001ux673a}

\passthrough{\lstinline!/bin/wifi!} 最终调用：

\begin{lstlisting}
wifi_connect(ssid, psk)
  -> SYS_wifi_connect
  -> sys_wifi_connect()
  -> mt7601u_connect(ssid, psk, 30000)
\end{lstlisting}

\passthrough{\lstinline!mt7601u\_connect()!} 在进程上下文执行 PBKDF2
4096 轮生成 32 字节 PMK，然后设置 目标 SSID，唤醒
\passthrough{\lstinline!mt7601u\_wq!} 的 delayed work，并睡眠等待最多 30
秒。

\begin{figure}[htbp]
\centering
\resizebox{\linewidth}{!}{\begin{tikzpicture}[node distance=6mm]
  \node[hbox] (scan) {SCAN\\channel 1--11};
  \node[box,right=of scan] (authp) {AUTH\_PENDING};
  \node[box,right=of authp] (auth) {AUTH\_SENT};
  \node[box,right=of auth] (assocp) {ASSOC\_PENDING};
  \node[box,right=of assocp] (assoc) {ASSOC\_SENT};
  \node[box,right=of assoc] (eapol) {ASSOCIATED\\EAPOL M1--M4};
  \node[hbox,right=of eapol] (key) {KEYED\\注册 \texttt{wlan0}};
  \draw[arr] (scan)--node[lab,above]{命中 SSID}(authp);
  \draw[arr] (authp)--(auth); \draw[arr] (auth)--(assocp); \draw[arr] (assocp)--(assoc);
  \draw[arr] (assoc)--(eapol); \draw[arr] (eapol)--(key);
  \draw[darr,bend right=35] (auth) to node[lab,below]{timeout} (authp);
  \draw[darr,bend right=38] (eapol) to node[lab,below]{watchdog} (assocp);
\end{tikzpicture}}
\caption{扫描、Open-System 认证、关联与 WPA2 状态机}
\label{fig:tiny-wifi-state}
\end{figure}


扫描在 2.4 GHz channel 1..11 间切换，每个 channel 解析 beacon/probe
response， 最多保存 32 个 AP。WPA2 四次握手中：

\begin{enumerate}
\def\labelenumi{\arabic{enumi}.}
\tightlist
\item
  收到 M1：保存 ANonce，生成 SNonce，PMK 派生 PTK，发送 M2；
\item
  收到 M3：校验 MIC，AES unwrap key data，提取 GTK；
\item
  把 PTK/GTK 写进 MT7601U hardware key tables，启用 AES-CCMP；
\item
  发送 M4，设置 \passthrough{\lstinline!handshake\_done!}。
\end{enumerate}

当前实现面向 WPA2-PSK/CCMP。没有完整支持
WPA3、802.1X/EAP、WEP、WPA-TKIP、 5 GHz、11n 高级速率控制或省电模式。

\section{RX：IRQ + NAPI
风格下半部}\label{rxirq--napi-ux98ceux683cux4e0bux534aux90e8}

probe 会在数据 bulk-IN endpoint 预提交 4 个 16 KiB buffer。数据到达时：

\begin{figure}[htbp]
\centering
\resizebox{\linewidth}{!}{\begin{tikzpicture}[node distance=5mm]
  \node[hbox] (mt) {MT7601U\\bulk-IN complete};
  \node[box,right=of mt] (irq) {EHCI IRQ55\\确认 qTD 状态};
  \node[box,right=of irq] (napi) {\texttt{mt7601u\_rx\_irq}\\调度 NAPI work};
  \node[box,right=of napi] (parse) {解析 RXWI / 802.11\\立即重提 RX buffer};
  \node[box,right=of parse] (kind) {管理帧/EAPOL 驱动内处理\\数据帧重建 Ethernet};
  \node[hbox,right=of kind] (net) {\texttt{net\_rx\_dev()}\\ARP/IP/UDP/TCP};
  \draw[arr] (mt)--(irq); \draw[arr] (irq)--(napi); \draw[arr] (napi)--(parse);
  \draw[arr] (parse)--(kind); \draw[arr] (kind)--(net);
\end{tikzpicture}}
\caption{Wi-Fi RX 从 EHCI 中断到 xv6 网络栈}
\label{fig:tiny-wifi-rx}
\end{figure}


硬中断只确认 EHCI 状态并调度 NAPI work，不在 IRQ 上下文解析
beacon、EAPOL 或 IP。 NAPI 每轮 budget 是 4；若一次用满
budget，workqueue 会继续调度。状态机的 timeout 使用独立
\passthrough{\lstinline!state\_work!}，因此 RX completion
与协议定时器互不阻塞。

\section{TX 与网络栈}\label{tx-ux4e0eux7f51ux7edcux6808}

网络栈交给 \passthrough{\lstinline!mt7601u\_net\_xmit()!} 的是 Ethernet
frame。驱动会：

\begin{enumerate}
\def\labelenumi{\arabic{enumi}.}
\tightlist
\item
  根据目标/源 MAC 和 EtherType 构造 802.11 ToDS header；
\item
  添加 LLC/SNAP；
\item
  关联后选择 AP 的 WCID；
\item
  握手完成后让硬件使用已安装的 CCMP key；
\item
  添加 MT7601U TXWI 和 USB DMA header；
\item
  通过 bulk-OUT 同步发送。
\end{enumerate}

RX 数据路径反向去掉 RXWI、802.11 header 和 LLC/SNAP，重建 Ethernet
header 后调用
\passthrough{\lstinline!net\_rx\_dev()!}。握手完成以前，普通数据 frame
不会交给 IP 层。

驱动注册的真实接口名是
\passthrough{\lstinline!wlan0!}。当前一条日志仍写成
\passthrough{\lstinline!wlan1 registered!}，这是日志
字符串遗留，不代表存在第二个接口；DHCP 应使用
\passthrough{\lstinline!wlan0!}。

\section{自动连接和手工测试}\label{ux81eaux52a8ux8fdeux63a5ux548cux624bux5de5ux6d4bux8bd5}

如果根文件系统有 \passthrough{\lstinline!/etc/wifi.conf!}，Tiny210
\passthrough{\lstinline!/init!} 会后台重试连接最多 10 次；
\passthrough{\lstinline!wifi!} 成功后执行：

\begin{lstlisting}[language=sh]
/bin/dhcp wlan0
\end{lstlisting}

手工连接：

\begin{lstlisting}[language=sh]
/bin/wifi YOUR_SSID YOUR_PASSPHRASE
/bin/dhcp wlan0
/bin/ping 192.168.0.1
\end{lstlisting}

也可直接让 \passthrough{\lstinline!wifi!} 读取配置：

\begin{lstlisting}[language=sh]
/bin/wifi
/bin/dhcp wlan0
\end{lstlisting}

成功的关键日志顺序通常是：

\begin{lstlisting}
ehci: high-speed device 148f:7601 ...
mt7601u: USB device 148f:7601 ASIC=...
firmware: mt7601u.bin loaded from /boot (... bytes)
mt7601u: MCU firmware running
mt7601u: permanent mac=...
mt7601u: MAC/BBP/RF bootstrap ready
mt7601u: RF calibration complete; bulk endpoints ready
mt7601u: target <ssid> found ... WPA2
mt7601u: Open-System authentication accepted
mt7601u: 802.11 associated ...
mt7601u: WPA2 four-way handshake complete
mt7601u: wlan1 registered ...     # 实际接口名 wlan0
dhcp: dev=wlan0 address=... gateway=... dns=...
\end{lstlisting}

\section{当前已知限制}\label{ux5f53ux524dux5df2ux77e5ux9650ux5236}

\begin{itemize}
\tightlist
\item
  只维护一个全局 \passthrough{\lstinline!mt7601u\_device!}，不支持多块
  MT7601U。
\item
  EHCI 启动只在约 2
  秒窗口内寻找设备，之后枚举线程永久睡眠；\textbf{不支持 USB 热插拔}。
  网卡应在启动前插好。
\item
  根口 full/low-speed 设备交给尚未实现的 OHCI；MT7601U 必须以 high-speed
  工作。
\item
  扫描只覆盖 channel 1..11，没有根据 country 自动扩大/缩小 regulatory
  domain。
\item
  \passthrough{\lstinline!mt7601u\_send\_assoc()!} 当前 association
  request 中仍有硬编码
  \passthrough{\lstinline!TP-Link\_B114!}，虽然扫描和
  \passthrough{\lstinline!wifi\_connect()!} 接受任意 SSID；连接其他 SSID
  前应把 这里改成
  \passthrough{\lstinline!dev->want\_ssid/dev->want\_ssid\_len!}。
\item
  TX bulk 是同步路径；高吞吐时会占用调用进程，尚无完整 TX
  queue/backpressure。
\item
  驱动没有 Linux mac80211/cfg80211 的速率控制、重连、扫描 API
  和监管框架。
\item
  D-cache 当前关闭，简化了 EHCI DMA 一致性；以后开启 D-cache 时必须补齐
  DMA coherent allocation 或 cache clean/invalidate。
\item
  Wi-Fi 配置是明文文件，系统没有密钥环或权限隔离。
\end{itemize}

\section{故障定位}\label{ux6545ux969cux5b9aux4f4d}

{\def\LTcaptype{none} % do not increment counter
\begin{longtable}[]{@{}p{0.34\linewidth}p{0.60\linewidth}@{}}
\toprule\noalign{}
最后日志/现象 & 优先检查 \\
\midrule\noalign{}
\endhead
\bottomrule\noalign{}
\endlastfoot
没有 \passthrough{\lstinline!initcall: ehci-s5p!} & 实板构建是否误用了
\passthrough{\lstinline!QEMU\_TEST=1!} \\
\passthrough{\lstinline!no high-speed device!} &
网卡是否开机前插入、USB-A 供电、PHY1/时钟、是否被识别为 full-speed \\
有 VID:PID，没有 probe & ID 是否在
\passthrough{\lstinline!mt7601u\_ids[]!}、usbcore/driver initcall
顺序 \\
\passthrough{\lstinline!firmware ... missing!} &
\passthrough{\lstinline!/boot/mt7601u.bin!} 和内置
\passthrough{\lstinline!firmware/mt7601u.bin!} \\
\passthrough{\lstinline!MCU firmware start timeout!} & bulk-OUT、FCE
地址/长度、固件版本、USB 错误 token \\
\passthrough{\lstinline!calibration ... failed!} & 第二 bulk-IN MCU
response endpoint、eFuse/RF、供电和 USB 稳定性 \\
能扫描但认证超时 & channel、BSSID、AP 是否允许 Open-System auth、USB RX
completion \\
关联后 M1 timeout & AP 安全模式应为 WPA2-PSK/CCMP，检查 EAPOL RX 和
receive filter \\
\passthrough{\lstinline!M3 MIC verification failed!} & 密码错误、PMK/PTK
派生或 replay/nonce \\
\passthrough{\lstinline!wifi: associated!} 但 DHCP 失败 & 使用
\passthrough{\lstinline!wlan0!}、ARP/DHCP response 是否进入
\passthrough{\lstinline!net\_rx\_dev()!} \\
一直连不上非 TP-Link\_B114 & 当前 association request 的硬编码 SSID
限制 \\
\end{longtable}
}
